\documentclass[11pt,a4paper]{article}

% --- Pakete für Mathematik und Symbole ---
\usepackage[utf8]{inputenc}
\usepackage[T1]{fontenc}
\usepackage[german]{babel}
\usepackage{amsmath, amssymb, amsthm, amsfonts}
\usepackage{geometry}
\geometry{a4paper, margin=2.5cm}
\usepackage{bm}
\usepackage{hyperref}
\usepackage{cite}

% --- Arithmetische Symbole ---
\newcommand{\Okt}{\mathbb{O}}
\newcommand{\Ham}{\mathbb{H}}
\newcommand{\Complex}{\mathbb{C}}
\newcommand{\Real}{\mathbb{R}}
\newcommand{\Nat}{\mathbb{N}}
\newcommand{\Zeta}{\zeta(s)}
\newcommand{\Gtwo}{G_2}
\newcommand{\SpinEight}{Spin(8)}

\title{Arithmetische Genetik: Die oktionische Triality-Helix und das spektrale Zerfallsmodell der natürlichen Zahlen}
\author{Thomas Hoffbauer}
\date{Februar 2026}

\begin{document}
	
	\maketitle
	
	\begin{abstract}
		Diese Arbeit stellt ein neuartiges mathematisches Framework vor, welches die natürliche Faktorisierung als energetischen Zerfallsprozess innerhalb einer oktionischen Hilbert-Raum-Struktur interpretiert. Basierend auf der Triality-Symmetrie der Gruppe $Spin(8)$ und der Einbettung in das exzeptionelle $E_8$-Gitter wird gezeigt, dass der Zahlenraum eine helikale Topologie aufweist. Wir führen die $\sqrt{18}$-Metrik als fundamentalen Kopplungsradius ein und korrelieren die spektrale Dynamik der Riemannschen Zeta-Funktion mit den Stabilitätsbedingungen des arithmetischen Genoms. Abschließend wird die Anwendung von Deep-Learning-Verfahren (AlphaFold-Adaption) zur Vorhersage arithmetischer Faltungsmuster diskutiert.
	\end{abstract}
	
	\section{Einleitung}
	Das traditionelle Verständnis der Primzahlverteilung als rein stochastischer Prozess wird hier durch ein deterministisches, geometrisches Modell ersetzt. Wir postulieren, dass die arithmetische Struktur der natürlichen Zahlen $n \in \Nat$ einer intrinsischen Genetik folgt, die auf der nicht-assoziativen Algebra der Oktaven (Oktionen) $\Okt$ basiert.
	
	\section{Topologie des arithmetischen Genoms}
	
	\subsection{Die Triality-Helix}
	Wir definieren eine Abbildung $\Phi: \Nat \to \Okt \times \Real$, wobei die vertikale Energieachse $E$ durch den Logarithmus der Zahl $t = \ln n$ gegeben ist. Die Triality-Symmetrie von $\SpinEight$ induziert eine zyklische Permutation der Primklassen (mod 12):
	\begin{equation}
		(E, A, B, C) \to (E, \text{S}_1, \text{S}_2, \text{S}_3)
	\end{equation}
	Diese Struktur bildet eine Doppelhelix, deren Pitch durch das Marion-Walter-Verhältnis von $1/10$ fixiert wird, um eine optimale Einbettung in den $G_2$-Invarianzraum zu gewährleisten.
	
	\subsection{Die $\sqrt{18}$-Metrik}
	Der fundamentale Abstand zwischen korrespondierenden Informationseinheiten auf der Helix ergibt sich aus der Kopplung der drei Triality-Sektoren. Innerhalb des $E_8$-Gitters (Kepler-Packing) wird die Stabilitätsmetrik $\delta_{G}$ definiert als:
	\begin{equation}
		\delta_{G} = \sqrt{3^2 + 3^2} = \sqrt{18}
	\end{equation}
	Diese Metrik schützt die arithmetische Information vor spektraler Dekohärenz.
	
	\section{Arithmetische Radioaktivität und Schwache Kraft}
	Zusammengesetzte Zahlen werden als metastabile Zustände modelliert. Die "schwache Kraft" vermittelt den Zerfall dieser Zustände in stabile primale Isotope. 
	\begin{itemize}
		\item \textbf{Stabile Zustände:} Primzahlen $p$ mit maximaler Bindungsenergie innerhalb der oktionischen Fano-Ebene.
		\item \textbf{Instabilitätsmaß:} Die Differenz zwischen dem Hecke-Eigenwert $\lambda_p$ und der Ramanujan-Schranke $2\sqrt{p}$ definiert das Zerfallspotenzial eines arithmetischen Nuklids.
	\end{itemize}
	
	\section{Zeta-Resonanz und DeepSearch-Integration}
	Die nichttrivialen Nullstellen $\gamma$ der Riemannschen Zeta-Funktion fungieren als die Resonanzfrequenzen des Systems. Ein harmonischer Zustand liegt vor, wenn:
	\begin{equation}
		\hat{D}_{DNA} = \sum_{\gamma} e^{i \gamma (\ln n + \phi_{Morley})} \to \text{min.}
	\end{equation}
	Durch den Einsatz von \textbf{DeepSearch-Architekturen} (Deep-Learning-basierte Proteinfaltung) können wir die energetisch günstigsten "Faltungen" von Primfaktor-Sequenzen vorhersagen. Dies erlaubt die Identifikation funktionaler arithmetischer Motive, die über rein statistische Methoden hinausgehen.
	
	\section{Fazit}
	Das vorgestellte Modell verknüpft die exzeptionelle Lie-Theorie mit der Biopolymertopologie. Es bietet eine konsistente Erklärung für die $\SpinEight$-Triality im Kontext der Zahlentheorie und liefert neue Impulse für die computergestützte Primzahlforschung.
	
\end{document}